% Living Showcase: Aussagen- und Hinweisboxen (Palette Slot 1).
% ZSFNewColumn: Kapitel bewusst auf neuer Spalte beginnen.
\ZSFNewColumn
\StartChapter[ch:boxes]{Aussagen und Hinweise}

\SubsectionBar[sec:box-core]{Titelboxen}

\begin{runintext}
  \ZSFkeyword{runintext} ist der ruhige Fliesstext-Baustein zwischen
  strukturierten Elementen. Er übernimmt automatisch die Lesbarkeitsregeln
  für schmale Spalten.
\end{runintext}

\begin{defbox}[Titelleiste und freier Inhalt\ZSFTitleTag{defbox}]
  \ZSFkeyword{defbox} trägt eine farbige Titelleiste und beliebigen Inhalt —
  gedacht für Definitionen und gewichtige Sätze.
  \ZSFhl{So sieht die hervorgehobene Kernaussage einer Box aus.}
\end{defbox}

\begin{defbox}[Dieselbe Box, leiser\ZSFTitleTag{weight=quiet}][weight=quiet]
  \ZSFkeyword*{weight=quiet} tauscht die Kopfzeile gegen einen dezenten linken
  Balken — für kompakte Aussagen. Kein zweiter Baustein, ein Regler.
\end{defbox}

\begin{warnbox}[Stolperfalle mit eigenem Block\ZSFTitleTag{warnbox}]
  \ZSFindexsee{Achtung}{warnbox}
  \ZSFkeyword{warnbox} ist der eigenständige Block für eine Warnung, die
  Raum braucht. Ohne Titelargument lautet die Leiste automatisch »Achtung«.
  \ZSFdanger{ZSFdanger ist die Inline-Pille.}
\end{warnbox}

\begin{defbox}[Dieselbe Box, anderer Ton\ZSFTitleTag{tone=neutral}][tone=neutral]
  Jede Box nimmt \ZSFkeyword*{tone} als zweites optionales Argument:
  \ZSFkeyword*{chapter} (Default), \ZSFkeyword*{neutral} für bewusst
  kapitelunabhängige Hinweise, \ZSFkeyword*{warn} für den Warn-Ton. Ein Regler
  statt einer neuen Box.
\end{defbox}

\begin{defbox}[Derselbe Regler, rahmenlos\ZSFTitleTag{tone=neutral}][weight=quiet, tone=neutral]
  Der Ton wirkt auch dort, wo es keine Kopfzeile gibt: Balken, Titelzeile und
  Aufzählungszeichen folgen ihm. Im Kapitelton bleibt die Fläche ungetönt,
  damit die Box dezent bleibt — erst ein gewählter Ton färbt sie.
\end{defbox}

\begin{defbox}[Regler greifen unabhängig\ZSFTitleTag{tone + frame}][tone=warn, frame=strong]
  \ZSFkeyword*{tone} wählt die Farbwelt, \ZSFkeyword*{frame} die Rahmenstärke
  — beide zusammen, in beliebiger Reihenfolge. Keine Wahl verwirft die andere.
\end{defbox}

\SubsubsectionBar[sec:box-hierarchy]{Dritte Lookup-Ebene}

\begin{defbox}[][weight=quiet]
  Der \ZSFkeyword{SubsubsectionBar} ist die optionale dritte Ebene, sichtbar
  kleiner und heller. Diese Box zeigt zugleich: ohne Titelargument entfällt
  die Kopfzeile ganz — dieselbe Box, ein Argument weniger.
\end{defbox}
